% ---------------------------------------------------------------------------
% Lista de abreviaturas (norma 3.11)
% ---------------------------------------------------------------------------
% Añade una fila por cada sigla que uses en el informe. Mantenla en orden
% alfabético. Elimina las que no llegues a utilizar.

\seccionSinNumerar{Lista de abreviaturas}

\begin{longtable}{@{}p{0.18\textwidth}p{0.78\textwidth}@{}}
\toprule
\textbf{Sigla} & \textbf{Significado} \\
\midrule
\endfirsthead
\toprule
\textbf{Sigla} & \textbf{Significado} \\
\midrule
\endhead
\bottomrule
\endfoot

ADC     & \textit{Analog-to-Digital Converter} — convertidor analógico-digital \\
ADCC    & \textit{ADC with Computation} — ADC con cómputo integrado \\
BM25    & \textit{Best Matching 25} — función de puntuación para recuperación léxica \\
BSR     & \textit{Bank Select Register} — registro de selección de banco \\
CCP     & \textit{Capture/Compare/PWM} \\
CLC     & \textit{Configurable Logic Cell} — celda lógica configurable \\
CVD     & \textit{Capacitive Voltage Divider} — divisor de tensión capacitivo \\
CWG     & \textit{Complementary Waveform Generator} \\
DAC     & \textit{Digital-to-Analog Converter} — convertidor digital-analógico \\
DFM     & \textit{Data Flash Memory} — memoria Flash de datos (EEPROM) \\
DPR     & \textit{Dense Passage Retrieval} — recuperación densa de pasajes \\
DSM     & \textit{Data Signal Modulator} — modulador de señal de datos \\
EEPROM  & \textit{Electrically Erasable Programmable Read-Only Memory} \\
EUSART  & \textit{Enhanced Universal Synchronous Asynchronous Receiver Transmitter} \\
FVR     & \textit{Fixed Voltage Reference} — referencia de tensión fija \\
GPIO    & \textit{General Purpose Input/Output} — entrada/salida de propósito general \\
GPR     & \textit{General Purpose Register} — registro de propósito general \\
HLT     & \textit{Hardware Limit Timer} \\
HLVD    & \textit{High/Low-Voltage Detect} — detección de alta y baja tensión \\
HNSW    & \textit{Hierarchical Navigable Small World} — grafo jerárquico para búsqueda aproximada \\
ICSP    & \textit{In-Circuit Serial Programming} \\
IDE     & \textit{Integrated Development Environment} — entorno de desarrollo integrado \\
I\textsuperscript{2}C & \textit{Inter-Integrated Circuit} \\
LLM     & \textit{Large Language Model} — modelo de lenguaje de gran escala \\
MCC     & \textit{MPLAB Code Configurator} \\
MCU     & \textit{Microcontroller Unit} — microcontrolador \\
MRR     & \textit{Mean Reciprocal Rank} — rango recíproco medio \\
MSSP    & \textit{Master Synchronous Serial Port} \\
NLP     & \textit{Natural Language Processing} — procesamiento de lenguaje natural \\
PFM     & \textit{Program Flash Memory} — memoria Flash de programa \\
PMD     & \textit{Peripheral Module Disable} — deshabilitación de módulos periféricos \\
PPS     & \textit{Peripheral Pin Select} — selección de pines de periférico \\
PWM     & \textit{Pulse-Width Modulation} — modulación por ancho de pulso \\
RAG     & \textit{Retrieval-Augmented Generation} — generación aumentada por recuperación \\
SFR     & \textit{Special Function Register} — registro de función especial \\
SPI     & \textit{Serial Peripheral Interface} \\
SRAM    & \textit{Static Random-Access Memory} \\
USB     & Universidad Simón Bolívar \\
WWDT    & \textit{Windowed Watchdog Timer} — temporizador de guardia con ventana \\
ZCD     & \textit{Zero-Cross Detection} — detección de cruce por cero \\
\end{longtable}
